\documentclass[10pt,a4paper]{amsart}
\usepackage{amsmath,amssymb,amsthm,graphicx}
% Retained preamble for the Stacks Project chapter morphisms.tex (Morphisms of Schemes).
% The chapter file pinned-original-source.tex carries only \input{preamble} on its first line, so no
% preamble text lives inside the chapter itself.  The Stacks Project uses one single shared preamble for
% every chapter; the snapshot stages it verbatim in sources/stacks-chow/preamble.tex.  The block below is
% copied mechanically, line for line, from that file's lines 14-18 (the Xy-pic setup) and lines 154-217
% (the native theorem declarations together with the mathematical macros).  Nothing was added, removed or
% redefined.  The wrapper class and the amsmath/amssymb/amsthm/graphicx base packages are supplied by the
% assembler, so the class/\externaldocument part of the shared preamble is not repeated here.
% For commutative diagrams we use Xy-pic
\usepackage[all]{xy}

% We use 2cell for 2-commutative diagrams.
\xyoption{2cell}
% Theorem environments.
%
\theoremstyle{plain}
\newtheorem{theorem}[subsection]{Theorem}
\newtheorem{proposition}[subsection]{Proposition}
\newtheorem{lemma}[subsection]{Lemma}

\theoremstyle{definition}
\newtheorem{definition}[subsection]{Definition}
\newtheorem{example}[subsection]{Example}
\newtheorem{exercise}[subsection]{Exercise}
\newtheorem{situation}[subsection]{Situation}

\theoremstyle{remark}
\newtheorem{remark}[subsection]{Remark}
\newtheorem{remarks}[subsection]{Remarks}

\numberwithin{equation}{subsection}

% Macros
%
\def\lim{\mathop{\mathrm{lim}}\nolimits}
\def\colim{\mathop{\mathrm{colim}}\nolimits}
\def\Spec{\mathop{\mathrm{Spec}}}
\def\Hom{\mathop{\mathrm{Hom}}\nolimits}
\def\Ext{\mathop{\mathrm{Ext}}\nolimits}
\def\SheafHom{\mathop{\mathcal{H}\!\mathit{om}}\nolimits}
\def\SheafExt{\mathop{\mathcal{E}\!\mathit{xt}}\nolimits}
\def\Sch{\mathit{Sch}}
\def\Mor{\mathop{\mathrm{Mor}}\nolimits}
\def\Ob{\mathop{\mathrm{Ob}}\nolimits}
\def\Sh{\mathop{\mathit{Sh}}\nolimits}
\def\NL{\mathop{N\!L}\nolimits}
\def\CH{\mathop{\mathrm{CH}}\nolimits}
\def\proetale{{pro\text{-}\acute{e}tale}}
\def\etale{{\acute{e}tale}}
\def\QCoh{\mathit{QCoh}}
\def\Ker{\mathop{\mathrm{Ker}}}
\def\Im{\mathop{\mathrm{Im}}}
\def\Coker{\mathop{\mathrm{Coker}}}
\def\Coim{\mathop{\mathrm{Coim}}}

% Boxtimes
%
\DeclareMathSymbol{\boxtimes}{\mathbin}{AMSa}{"02}

%
% Macros for moduli stacks/spaces
%
\def\QCohstack{\mathcal{QC}\!\mathit{oh}}
\def\Cohstack{\mathcal{C}\!\mathit{oh}}
\def\Spacesstack{\mathcal{S}\!\mathit{paces}}
\def\Quotfunctor{\mathrm{Quot}}
\def\Hilbfunctor{\mathrm{Hilb}}
\def\Curvesstack{\mathcal{C}\!\mathit{urves}}
\def\Polarizedstack{\mathcal{P}\!\mathit{olarized}}
\def\Complexesstack{\mathcal{C}\!\mathit{omplexes}}
% \Pic is the operator that assigns to X its picard group, usage \Pic(X)
% \Picardstack_{X/B} denotes the Picard stack of X over B
% \Picardfunctor_{X/B} denotes the Picard functor of X over B
\def\Pic{\mathop{\mathrm{Pic}}\nolimits}
\def\Picardstack{\mathcal{P}\!\mathit{ic}}
\def\Picardfunctor{\mathrm{Pic}}
\def\Deformationcategory{\mathcal{D}\!\mathit{ef}}

% --- appended 2026-10-12 (Chief): the project own definitions of the reference/slogan/history
% display environments, copied mechanically from the shared Stacks preamble lines 8-12
% (\usepackage{verbatim} plus the three \newenvironment lines, which the chapter preamble keeps
% verbatim from the project preamble). Nothing added, removed or redefined: these environments
% carry no mathematical content and the project itself typesets them as comments. Without these
% three lines a counted statement whose own text uses them fails with Environment slogan undefined.
\usepackage{verbatim}
\newenvironment{reference}{\comment}{\endcomment}
\newenvironment{slogan}{\comment}{\endcomment}
\newenvironment{history}{\comment}{\endcomment}
\usepackage{cleveref}
\crefname{originallemma}{original lemma}{original lemmas}
\Crefname{originallemma}{Original lemma}{Original lemmas}
\makeatletter
\newlabel{constructions-lemma-proj-separated}{{01MC}{source}{}{sourcepointer-1}{}}
\newlabel{constructions-lemma-proj-separated@cref}{{[originallemma][0][]01MC}{[1][0][]source}{}{}{}}
\makeatother
\crefname{originallemma}{original lemma}{original lemmas}
\Crefname{originallemma}{Original lemma}{Original lemmas}
\makeatletter
\newlabel{constructions-lemma-twisting-and-proj}{{02NC}{source}{}{sourcepointer-2}{}}
\newlabel{constructions-lemma-twisting-and-proj@cref}{{[originallemma][0][]02NC}{[1][0][]source}{}{}{}}
\makeatother
\crefname{originallemma}{original lemma}{original lemmas}
\Crefname{originallemma}{Original lemma}{Original lemmas}
\makeatletter
\newlabel{lemma-quasi-projective-open-projective}{{07RM}{source}{}{sourcepointer-3}{}}
\newlabel{lemma-quasi-projective-open-projective@cref}{{[originallemma][0][]07RM}{[1][0][]source}{}{}{}}
\makeatother
\crefname{originaldefinition}{original definition}{original definitions}
\Crefname{originaldefinition}{Original definition}{Original definitions}
\makeatletter
\newlabel{properties-definition-ample}{{01PS}{source}{}{sourcepointer-4}{}}
\newlabel{properties-definition-ample@cref}{{[originaldefinition][0][]01PS}{[1][0][]source}{}{}{}}
\makeatother
\crefname{originallemma}{original lemma}{original lemmas}
\Crefname{originallemma}{Original lemma}{Original lemmas}
\makeatletter
\newlabel{properties-lemma-ample-immersion-into-proj}{{01Q1}{source}{}{sourcepointer-5}{}}
\newlabel{properties-lemma-ample-immersion-into-proj@cref}{{[originallemma][0][]01Q1}{[1][0][]source}{}{}{}}
\makeatother
\crefname{originalproposition}{original proposition}{original propositions}
\Crefname{originalproposition}{Original proposition}{Original propositions}
\makeatletter
\newlabel{properties-proposition-characterize-ample}{{01Q3}{source}{}{sourcepointer-6}{}}
\newlabel{properties-proposition-characterize-ample@cref}{{[originalproposition][0][]01Q3}{[1][0][]source}{}{}{}}
\makeatother
\title{Over a base carrying an ample invertible sheaf, projective morphisms are H-projective and quasi-projective ones are H-quasi-projective}
\author{The Stacks Project Authors}
\date{stacks-project:morphisms.tex:local-snapshot-20261012; GNU FDL 1.2 or later}
\begin{document}
\maketitle
\noindent\textit{Editorial note.} All delivered mathematical text is the author's own, copied from the pinned Stacks Project chapter file pinned-original-source.tex; no statement or proof step was rewritten, reordered or removed.\par\medskip
\noindent\textit{Editorial note.} Following the multi-claim convention for this source, source\_url is the source-level URL https://stacks.math.columbia.edu/ and the source label lemma-projective-over-quasi-projective-is-H-projective is carried in claim\_key and here rather than in a per-claim tag URL.\par\medskip
\noindent\textit{Editorial note.} The preamble is the chapter's shared Stacks Project preamble, retained as sources/stacks-morphisms/retained-preamble.tex: a line-for-line excerpt of the shared Stacks preamble staged with this snapshot. No macro was invented or redefined, and this card adds no preamble text of its own.\par\medskip
\noindent\textit{Editorial note.} Statement span 10847-10854 and proof span 10856-10882 come from the halfA pool card operations/pool-cards/stacks-morphisms-halfA.json and were re-read in the pinned chapter file before this card was written.\par\medskip
\noindent\textit{Editorial note.} The proof uses the author's identification of a projective bundle of a free module with a projective space over the base and the invariance of projective bundles under twisting; both are preserved as written.\par\medskip
\noindent\textit{Editorial note.} 6 cross-references occur inside the counted spans; they are listed in external\_references as navigation pointers only.\par\medskip
\noindent\textit{Allowed standard prerequisites (editorial).}\begin{itemize}
\item A scheme carrying an ample invertible sheaf is quasi-compact and separated, and it admits an immersion into a projective space over the integers, cited from the properties and constructions chapters. Cited at proof lines 10856-10860. Supplies the hypotheses under which the passage from H-projective to H-quasi-projective is available.
\item The statement that over a base with those properties a quasi-projective morphism is obtained from an H-projective one, cited from the same chapter. Cited at proof lines 10859-10861. Gives that the first assertion implies the second.
\item The construction of a morphism with a projective bundle as source from a generating set of sections of the corresponding module, the identification of a projective bundle of a free module with a projective space over the base, and the invariance of a projective bundle under twisting the module by an invertible sheaf. Cited at proof lines 10868-10881. Produces the closed immersion into a projective space over the base in the generated case, and reduces the general case to it by an ample twist.
\end{itemize}\medskip
\ifdefined\hypertarget\hypertarget{sourcepointer-1}{}\fi
\noindent\textit{Source reference.} Original lemma 01MC: focused-seeds/AGP/sources/stacks-snapshot-20261011/tags.txt, source lines 2117--2117. Same-chapter reference used in the retained author text at source line(s) 9142, 10597, 10860. It is supplied only as a navigation pointer so that the retained text compiles; the counted item delivers the author's own statement and the author's complete proof, and no part of the referenced item is copied in or substituted for the counted argument.\par\medskip
\ifdefined\hypertarget\hypertarget{sourcepointer-2}{}\fi
\noindent\textit{Source reference.} Original lemma 02NC: focused-seeds/AGP/sources/stacks-snapshot-20261011/tags.txt, source lines 3475--3475. Same-chapter reference used in the retained author text at source line(s) 10880. It is supplied only as a navigation pointer so that the retained text compiles; the counted item delivers the author's own statement and the author's complete proof, and no part of the referenced item is copied in or substituted for the counted argument.\par\medskip
\ifdefined\hypertarget\hypertarget{sourcepointer-3}{}\fi
\noindent\textit{Source reference.} Original lemma 07RM: focused-seeds/AGP/sources/stacks-snapshot-20261011/tags.txt, source lines 9872--9872. Same-chapter reference used in the retained author text at source line(s) 10508, 10792, 10861. It is supplied only as a navigation pointer so that the retained text compiles; the counted item delivers the author's own statement and the author's complete proof, and no part of the referenced item is copied in or substituted for the counted argument.\par\medskip
\ifdefined\hypertarget\hypertarget{sourcepointer-4}{}\fi
\noindent\textit{Source reference.} Original definition 01PS: focused-seeds/AGP/sources/stacks-snapshot-20261011/tags.txt, source lines 2241--2241. Same-chapter reference used in the retained author text at source line(s) 8623, 8846, 10858. It is supplied only as a navigation pointer so that the retained text compiles; the counted item delivers the author's own statement and the author's complete proof, and no part of the referenced item is copied in or substituted for the counted argument.\par\medskip
\ifdefined\hypertarget\hypertarget{sourcepointer-5}{}\fi
\noindent\textit{Source reference.} Original lemma 01Q1: focused-seeds/AGP/sources/stacks-snapshot-20261011/tags.txt, source lines 2250--2250. Same-chapter reference used in the retained author text at source line(s) 8742, 10859, 10909. It is supplied only as a navigation pointer so that the retained text compiles; the counted item delivers the author's own statement and the author's complete proof, and no part of the referenced item is copied in or substituted for the counted argument.\par\medskip
\ifdefined\hypertarget\hypertarget{sourcepointer-6}{}\fi
\noindent\textit{Source reference.} Original proposition 01Q3: focused-seeds/AGP/sources/stacks-snapshot-20261011/tags.txt, source lines 2252--2252. Same-chapter reference used in the retained author text at source line(s) 8676, 8876, 9431, 9484, 9655, 10874. It is supplied only as a navigation pointer so that the retained text compiles; the counted item delivers the author's own statement and the author's complete proof, and no part of the referenced item is copied in or substituted for the counted argument.\par\medskip
\begin{lemma}
\label{lemma-projective-over-quasi-projective-is-H-projective}
Let $S$ be a scheme which admits an ample invertible sheaf. Then
\begin{enumerate}
\item any projective morphism $X \to S$ is H-projective, and
\item any quasi-projective morphism $X \to S$ is H-quasi-projective.
\end{enumerate}
\end{lemma}

\begin{proof}
The assumptions on $S$ imply that $S$ is quasi-compact and separated, see
Properties, Definition \ref{properties-definition-ample} and
Lemma \ref{properties-lemma-ample-immersion-into-proj}
and Constructions, Lemma \ref{constructions-lemma-proj-separated}.
Hence Lemma \ref{lemma-quasi-projective-open-projective}
applies and we see that (1) implies (2).
Let $\mathcal{E}$ be a finite type quasi-coherent $\mathcal{O}_S$-module.
By our definition of projective morphisms it suffices to show that
$\mathbf{P}(\mathcal{E}) \to S$ is H-projective.
If $\mathcal{E}$ is generated by finitely many global sections,
then the corresponding surjection $\mathcal{O}_S^{\oplus n} \to \mathcal{E}$
induces a closed immersion
$$
\mathbf{P}(\mathcal{E}) \longrightarrow
\mathbf{P}(\mathcal{O}_S^{\oplus n}) = \mathbf{P}^{n - 1}_S
$$
as desired. In general, let $\mathcal{L}$ be an ample invertible sheaf on $S$.
By Properties, Proposition \ref{properties-proposition-characterize-ample}
there exists an integer $n$ such that
$\mathcal{E} \otimes_{\mathcal{O}_S} \mathcal{L}^{\otimes n}$
is globally generated by finitely many sections. Since
$\mathbf{P}(\mathcal{E}) =
\mathbf{P}(\mathcal{E} \otimes_{\mathcal{O}_S} \mathcal{L}^{\otimes n})$ by
Constructions, Lemma \ref{constructions-lemma-twisting-and-proj}
this finishes the proof.
\end{proof}

\end{document}
